\Needspace{6\baselineskip}
\begin{center}\hypertarget{native-257-article184}{}\pdfbookmark[2]{184.}{article-184}184.\end{center}
Ut disquisitiones praecedentes per exemplum illustrentur, supponamus,
per observationes, in quibus praecisio aequalis praesumenda sit, inventum esse
\[
\begin{aligned}
p-q+2r&=3\\
3p+2q-5r&=5\\
4p+q+4r&=21,
\end{aligned}
\]
\SourceLeaf{258}per quartam vero, cui praecisio dimidia tantum tribuenda est, prodiisse
\[
-2p+6q+6r=28.
\]
Loco aequationis ultimae itaque hanc substituemus
\[
-p+3q+3r=14
\]
hancque ex observatione prioribus praecisione aequali provenisse supponemus.
Hinc fit
\[
\begin{aligned}
P&=27p+6q-88\\
Q&=6p+15q+r-70\\
R&=q+54r-107,
\end{aligned}
\]
atque hinc per eliminationem
\[
\begin{aligned}
19899p&=49154+809P-324Q+6R\\
737q&=2617-12P+54Q-R\\
6633r&=12707+2P-9Q+123R.
\end{aligned}
\]
Incognitarum itaque valores maxime probabiles erunt
\[
\begin{aligned}
p&=2{,}470\\
q&=3{,}551\\
r&=1{,}916,
\end{aligned}
\]
atque praecisio relativa his determinationibus tribuenda, posita praecisione
observationum primitivarum $=1$,
\[
\begin{aligned}
\text{pro }p\quad .\,.\,.\,.\quad &\BareRoot{\frac{19899}{809}}=4{,}96\\[4pt]
\text{pro }q\quad .\,.\,.\,.\quad &\BareRoot{\frac{737}{54}}=3{,}69\\[4pt]
\text{pro }r\quad .\,.\,.\,.\quad &\BareRoot{\frac{2211}{41}}=7{,}34.
\end{aligned}
\]
